

 \LabHeader{5}{1}

\begin{enumerate}
\item \textbf{Всі прості.} Дано натуральні числа $a$ і $b$. Вивести на екран усі прості числа з проміжку $[a; b]$ і знайти їхню суму.
\item \textbf{Сума.} Дано натуральне число $n$. Знайти значення виразу
$$
\frac{\cos 1}{\sin 1}\cdot\frac{\cos 1+\cos 2}{\sin 1+\sin 2}\cdot\ldots\cdot\frac{\cos 1+\cos 2 +\ldots+\cos n}{\sin 1+\sin 2+\ldots+\sin n}
$$
з точністю до чотирьох знаків після коми.
\item \textbf{Табулювання функції.} Вивести у вигляді таблиці значення функції
$$
f(x)=x^2-2x+9
$$
для $x \in [0; 2]$ з кроком $ h = 0.2 $. Знайти значення аргументу, за якого функція досягає найменшого значення. Вивести повідомлення виду

{\ttfamily\color{blue} Мінімальне значення функції дорівнює $\ldots$  при  $x = \ldots $ }
\end{enumerate}

 

 \LabHeader{5}{2}

\begin{enumerate}
\item \textbf{Досконалі числа.} Натуральне число називається {\itshape досконалим}, якщо воно дорівнює
сумі своїх дільників, відмінних від самого числа. Наприклад, $6 =
1+2+3$. Знайти всі досконалі числа від 1 до 1000000.
\item \textbf{Послідовність.} Користувач вводить з клавіатури натуральне число $N$, а потім по\-слі\-дов\-ність з $N$ цілих чисел. Чи правда, що серед них немає від'ємних?
\item \textbf{Табулювання функції.} Вивести у вигляді таблиці значення функції
$$
f(x)=-2x^2-2x+1
$$
для $ x \in [-1; 0]$ з кроком $h=0.1$. Знайти значення аргументу, за якого функція досягає найбільшого значення. Вивести повідомлення виду

 {\ttfamily\color{blue} Максимальне значення функції дорівнює $\ldots$  при  $x = \ldots $ }
\end{enumerate}

 

 \LabHeader{5}{3}

\begin{enumerate}
\item \textbf{Числа Армстронга.} У теорії чисел прийнято таке визначення: $m$-цифрове натуральне число називається {\itshape самозакоханим} (або {\itshape числом Армстронга}), якщо воно дорівнює сумі $m$-х степенів своїх цифр. Наприклад, при $m=3$ маємо $153=1^3+5^3+3^3$. А чи існують інші самозакохані трицифрові числа?
\item \textbf{Добуток.} Знайти добуток
$$
P=\frac{2}{1}\cdot \frac{2}{3}\cdot \frac{4}{3}\cdot \frac{4}{5}\cdot \frac{6}{5} \cdot \frac{6}{7}\cdot\ldots\cdot\frac{1000}{999}\cdot\frac{1000}{1001}.
$$
\item \textbf{Табулювання функції.} Вивести у вигляді таблиці значення функції
$$
f(x)=18x^2-9x-123
$$
для $ x \in [0; 2] $ з кроком $ h = 0.25 $. Знайти значення аргументу, за якого функція досягає найменшого значення. Вивести повідомлення виду

{\ttfamily\color{blue} Мінімальне значення функції дорівнює $\ldots$ при $x = \ldots$}
\end{enumerate}

 

 \LabHeader{5}{4}

\begin{enumerate}
\item \textbf{Найбільший непарний дільник.} Дано натуральне число $n$. Знайти його найбільший непарний дільник, відмінний від самого числа.
\item \textbf{Корені.} Знайти значення виразу
$$
S=\underbrace{\sqrt{2+\sqrt{2+\sqrt{2+\ldots}}}}_{\text{усього 100 знаків радикала}}
$$
з точністю до чотирьох знаків після коми.
\item \textbf{Табулювання функції.} Вивести у вигляді таблиці значення функції
$$
f(x)=-23x^2+46x+97
$$
для $x \in [0; 2]$ з кроком $ h = 0.2 $. Знайти значення аргументу, за якого функція досягає найбільшого значення. Вивести повідомлення виду

 {\ttfamily\color{blue} Максимальне значення функції дорівнює $\ldots$  при  $x = \ldots $ }
\end{enumerate}


 

 \LabHeader{5}{5}

\begin{enumerate}
\item \textbf{Двоцифрові числа.} Знайти суму всіх двоцифрових натуральних чисел, перша цифра яких більша за другу.
\item \textbf{Послідовність.} Користувач вводить з клавіатури натуральне число $N$, а потім по\-слі\-до\-вність з $N$ цілих чисел. Чи правда, що серед них є хоча б одне від'ємне?
\item \textbf{Табулювання функції.} Вивести у вигляді таблиці значення функції
$$
f(x)=15x^2+10x-1
$$
для $ x \in [-1; 0]$ з кроком $h=0.1$. Знайти значення аргументу, за якого функція досягає найменшого значення. Вивести повідомлення виду

 {\ttfamily\color{blue} Мінімальне значення функції дорівнює $\ldots$  при  $x = \ldots $ }
\end{enumerate}

 

 \LabHeader{5}{6}

\begin{enumerate}
\item \textbf{Простий дільник.} Дано складене натуральне число $n$. Знайти його найбільший простий дільник.
\item \textbf{Синуси.} Дано натуральне число $n$. Обчислити суму
$$
S=\frac{1}{\sin 1}+\frac{1}{\sin 1+\sin 2}+\ldots+\frac{1}{\sin 1+\sin 2+\ldots+\sin n}
$$
з точністю до чотирьох знаків після коми.
\item \textbf{Табулювання функції.} Вивести у вигляді таблиці значення функції
$$
f(x)=-2x^2+5x-11
$$
для $x \in [0; 2]$ з кроком $ h = 0.25 $. Знайти значення аргументу, за якого функція досягає найбільшого значення. Вивести повідомлення виду

 {\ttfamily\color{blue} Максимальне значення функції дорівнює $\ldots$  при  $x = \ldots $ }
\end{enumerate}


 

 \LabHeader{5}{7}

\begin{enumerate}
\item \textbf{Трицифрові числа.} Знайти всі трицифрові числа, складені з трьох різних цифр, які діляться на квадрат суми своїх цифр.
\item \textbf{Синуси.} Дано натуральне число $n$ і дійсне $x$. Знайти значення виразу
$$
S=\sin x+\sin x^2 +\ldots +\sin x^n
$$
з точністю до чотирьох знаків після коми.

{\itshape Намагайтеся розв'язати це завдання без використання функції піднесення до степеня.}
\item \textbf{Табулювання функції.} Вивести у вигляді таблиці значення функції
$$
f(x)=x^2+20x-1
$$
для $ x \in [-10; -9]$ з кроком $h=0.1$. Знайти значення аргументу, за якого функція досягає найменшого значення. Вивести повідомлення виду

 {\ttfamily\color{blue} Мінімальне значення функції дорівнює $\ldots$  при  $x = \ldots $ }
\end{enumerate}


 

 \LabHeader{5}{8}

\begin{enumerate}
\item \textbf{Трицифрові числа.} Знайти суму всіх трицифрових натуральних чисел, перша цифра яких дорівнює сумі двох інших.
\item \textbf{Послідовність.} Користувач вводить з клавіатури натуральне число $N$, а потім по\-слі\-до\-вність з $N$ цілих чисел. Знайти середнє арифметичне непарних додатних елементів. Якщо таких елементів не було, вивести відповідне повідомлення.
\item \textbf{Табулювання функції.} Вивести у вигляді таблиці значення функції
$$
f(x)=-x^2+5x+10
$$
для $ x \in [2; 3] $ з кроком $ h = 0.1 $. Знайти значення аргументу, за якого функція досягає найбільшого значення. Вивести повідомлення виду

 {\ttfamily\color{blue} Максимальне значення функції дорівнює $\ldots$  при  $x = \ldots $ }
\end{enumerate}


 

 \LabHeader{5}{9}

\begin{enumerate}
\item \textbf{Двоцифрові числа.} Знайти суму всіх двоцифрових натуральних чисел, друга цифра яких більша за першу.
\item \textbf{Добуток.} Дано натуральне число $n>2$. Знайти добуток
$$
P=\left(1-\frac{1}{2^2}\right)\left(1-\frac{1}{3^2}\right)\cdot\ldots\cdot\left(1-\frac{1}{n^2}\right)
$$
з точністю до чотирьох знаків після коми.
\item \textbf{Табулювання функції.} Вивести у вигляді таблиці значення функції
$$
f(x)=2x^2+8x-10
$$
для $ x \in [-3.5; -1]$ з кроком $h=0.25$. Знайти значення аргументу, за якого функція досягає найменшого значення. Вивести повідомлення виду

 {\ttfamily\color{blue} Мінімальне значення функції дорівнює $\ldots$  при  $x = \ldots $ }
\end{enumerate}


 

 \LabHeader{5}{10}

\begin{enumerate}
\item \textbf{Числа Армстронга.} У теорії чисел прийнято таке визначення: $m$-цифрове натуральне число називається {\itshape самозакоханим} (або {\itshape числом Армстронга}), якщо воно дорівнює сумі $m$-х степенів своїх цифр. Наприклад, при $m=3$ маємо $153=1^3+5^3+3^3$. А чи існують інші самозакохані трицифрові числа?
\item \textbf{Сума.} Дано натуральне число $n$. Знайти суму перших $n$ доданків ряду
$$
\frac{2}{1}+\frac{3}{2}+\frac{4}{3}+\ldots
$$
з точністю до чотирьох знаків після коми.
\item \textbf{Табулювання функції.} Вивести у вигляді таблиці значення функції
$$
f(x)=-x^2-3x+110
$$
для $ x \in [-2; 0]$ з кроком $h=0.1$. Знайти значення аргументу, за якого функція досягає найбільшого значення. Вивести повідомлення виду

 {\ttfamily\color{blue} Максимальне значення функції дорівнює $\ldots$  при  $x = \ldots $ }
\end{enumerate}


 

 \LabHeader{5}{11}

\begin{enumerate}
\item \textbf{Трицифрові числа.} Вивести всі трицифрові натуральні числа, сума цифр яких дорівнює 11, а також їх кількість.
\item \textbf{Сума.} Знайти суму
$$
S=-1^2+2^2-3^2+4^2-\dots+100^2.
$$
\item \textbf{Табулювання функції.} Вивести у вигляді таблиці значення функції
$$
f(x)=2x^2-9x-10
$$
для $ x \in [1; 3] $ з кроком $ h = 0.25 $. Знайти значення аргументу, за якого функція досягає найменшого значення. Вивести повідомлення виду

 {\ttfamily\color{blue} Мінімальне значення функції дорівнює $\ldots$  при  $x = \ldots $ }
\end{enumerate}


 

 \LabHeader{5}{12}

\begin{enumerate}
\item \textbf{Двоцифрові числа.} Вивести в рядок через пробіл усі двоцифрові натуральні числа, сума цифр яких кратна трьом. Знайти суму всіх таких чисел.
\item \textbf{Сума.} Обчислити суму
$$
S=\frac{1}{2}-\frac{2}{3}+\frac{3}{4}-\frac{4}{5}+\ldots+\frac{111}{112}
$$
з точністю до чотирьох знаків після коми.
\item \textbf{Табулювання функції.} Вивести у вигляді таблиці значення функції
$$
f(x)=-x^2+3x+56
$$
для $ x \in [0; 2] $ з кроком $ h = 0.1 $. Знайти значення аргументу, за якого функція досягає найбільшого значення. Вивести повідомлення виду

 {\ttfamily\color{blue} Максимальне значення функції дорівнює $\ldots$  при  $x = \ldots $ }
\end{enumerate}


 

 \LabHeader{5}{13}

\begin{enumerate}
\item \textbf{Паліндроми.} Знайти суму всіх трицифрових натуральних паліндромів.
\item \textbf{Сума.} Дано натуральне число $n$. Знайти суму перших $n$ доданків ряду
$$
1+\frac{1}{\sqrt{3}}+\frac{1}{\sqrt{5}}+\frac{1}{\sqrt{7}}+\ldots
$$
з точністю до чотирьох знаків після коми.
\item \textbf{Табулювання функції.} Вивести у вигляді таблиці значення функції
$$
f(x)=2x^2-15x-8
$$
для $ x \in [2; 4] $ з кроком $ h = 0.25 $. Знайти значення аргументу, за якого функція досягає найменшого значення. Вивести повідомлення виду

 {\ttfamily\color{blue} Мінімальне значення функції дорівнює $\ldots$  при  $x = \ldots $ }
\end{enumerate}


 

 \LabHeader{5}{14}

\begin{enumerate}
\item \textbf{Тільки непарні.} Дано цілі числа $a$ і $b$. Вивести на екран через пробіл усі непарні числа з проміжку $[a; b]$ і знайти їхню суму.
\item \textbf{Формула Рамануджана}\footnote{Срініваза Рамануджан (1887--1920) "--- індійський математик-самоучка, першовідкривач багатьох дивовижних формул. Також відомий своїми роботами в~галузі теорії чисел.}.
$$
\footnotesize
\sqrt{1\mathstrut+2\sqrt{1\mathstrut+3\sqrt{1\mathstrut+4\sqrt{1\mathstrut+\ldots 99\sqrt{1\mathstrut+100}}}}} \quad \text {(усього 99 знаків радикала).}
$$
\item \textbf{Табулювання функції.} Вивести у вигляді таблиці значення функції
$$
f(x)=-2x^2-3x+1
$$
для $ x \in [-2; 0]$ з кроком $h=0.25$. Знайти значення аргументу, за якого функція досягає найбільшого значення. Вивести повідомлення виду

 {\ttfamily\color{blue} Максимальне значення функції дорівнює $\ldots$  при  $x = \ldots $ }
\end{enumerate}


 

 \LabHeader{5}{15}

\begin{enumerate}
\item \textbf{Трицифрові числа.} Знайти суму всіх трицифрових натуральних чисел, третя цифра яких дорівнює сумі двох інших.
\item \textbf{Найбільша відстань.} Користувач вводить з клавіатури натуральне число $N$, а потім по\-слі\-до\-вність з $N$ довільних чисел. Знайти пару елементів цієї послідовності, відстань між якими на числовій осі найбільша.
\item \textbf{Табулювання функції.} Вивести у вигляді таблиці значення функції
$$
f(x)=x^2+15x+8
$$
для $ x \in [-10; -7] $ з кроком $ h = 0.25 $. Знайти значення аргументу, за якого функція досягає найменшого значення. Вивести повідомлення виду

 {\ttfamily\color{blue} Мінімальне значення функції дорівнює $\ldots$  при  $x = \ldots $ }
\end{enumerate}



